\section{Analysis}

The results clarify what each layer measures. L1 reconstruction identifies whether an explanation contains enough information to recover the observed edit, but it is highly sensitive to answer copying. This is useful for detecting leakage, not for claiming internal causal faithfulness.

L2 is the more behaviorally meaningful layer, but the pilot also shows why it is difficult. Rule-relevant variants often create competing model edits rather than a clean cancel/change label. We therefore keep competing edits as a separate class and avoid binary preserve-versus-cancel scoring.

L3 rule/evidence checks are useful for rejecting generic or mismatched explanations, yet the current verifier is lexical and benefits from automatic rule templates. Its 0.767 Macro-F1 should be interpreted as evidence that the constructed benchmark is separable, not as evidence that the verifier understands grammar.

Finally, the reranking application demonstrates why evaluation must include reward-hacking checks. The best automatic reranker prefers explicit templates that repeat the edit. This is acceptable as a diagnostic finding, but not as a learner-facing explanation-selection result.

The Qwen3 method pilot shows the same lesson on natural teacher candidates. Automatic gates are useful for eliminating obvious failures, but the final Codex/AI pseudo-validation still rejects 11 of 41 candidates that had passed automatic repair and revalidation. Several failures involve linguistically implausible rules rather than edit-alignment errors. This supports the design choice to keep rule and evidence verification separate from reconstruction, and it also prevents the current paper from overclaiming that automatic RuleFaith gates alone solve explanation faithfulness.
